% !PRINTED 582
\weberpage{582}
%<<CONT
und dass sie also, da sie nur vom $(n-1)^{\text{ten}}$ Grade ist, identisch $=1$ sein
muss. Daraus folgt nach (14):
\[
1 = t^{n-1} S\,\frac{F_0(\tau)}{F'(\tau)} + t^{n-2} S\,\frac{F_1(\tau)}{F'(\tau)} + \cdots + S\,\frac{F_{n-1}(\tau)}{F'(\tau)}
\]
oder
\weq{17}{\begin{aligned} S\,\frac{F_\nu(\tau)}{F'(\tau)} &= 0, \quad S\,\frac{F_{n-1}(\tau)}{F'(\tau)} = 1,\\ \nu &= 0, 1, \ldots, n-2. \end{aligned}}
Hieraus ergiebt sich mit Benutzung von (16):
\weq{18}{S\,\frac{\omega}{F'(\tau)} = B_{n-1},}
oder der Satz:

4. \gesperrt{Ist $\omega$ eine ganze Zahl oder ein ganzes Functional des K\"orpers $\Omega$, so ist die Spur von $\dfrac{\omega}{F'(\tau)}$ ein ganzes rationales Functional.}

\begin{center}
\rule{0.2\linewidth}{0.4pt}
\end{center}
